\documentclass[10pt,a4paper]{amsart}
\usepackage{amsmath,amssymb,amsthm,graphicx}
% Retained minimal preamble for the Stacks Project chapter more-algebra.tex (More on Algebra).
% The chapter file pinned-original-source.tex carries only \input{preamble} on its first line, so no
% preamble text lives inside the chapter itself.  The Stacks Project uses one single shared preamble for
% every chapter; the snapshot stages it identically as sources/stacks-chow/preamble.tex and
% sources/stacks-etale-cohomology/preamble.tex.  The block below is copied mechanically, line for line,
% from sources/stacks-chow/preamble.tex lines 14-18 (the Xy-pic setup) and 154-217 (the native theorem
% declarations together with the mathematical macros).  Nothing was added, removed or redefined; the
% wrapper class and the amsmath/amssymb/amsthm/graphicx base packages are supplied by the assembler.
% Chapter: more-algebra.tex -- More on Algebra.
% For commutative diagrams we use Xy-pic
\usepackage[all]{xy}

% We use 2cell for 2-commutative diagrams.
\xyoption{2cell}

% Theorem environments.
%
\theoremstyle{plain}
\newtheorem{theorem}[subsection]{Theorem}
\newtheorem{proposition}[subsection]{Proposition}
\newtheorem{lemma}[subsection]{Lemma}

\theoremstyle{definition}
\newtheorem{definition}[subsection]{Definition}
\newtheorem{example}[subsection]{Example}
\newtheorem{exercise}[subsection]{Exercise}
\newtheorem{situation}[subsection]{Situation}

\theoremstyle{remark}
\newtheorem{remark}[subsection]{Remark}
\newtheorem{remarks}[subsection]{Remarks}

\numberwithin{equation}{subsection}

% Macros
%
\def\lim{\mathop{\mathrm{lim}}\nolimits}
\def\colim{\mathop{\mathrm{colim}}\nolimits}
\def\Spec{\mathop{\mathrm{Spec}}}
\def\Hom{\mathop{\mathrm{Hom}}\nolimits}
\def\Ext{\mathop{\mathrm{Ext}}\nolimits}
\def\SheafHom{\mathop{\mathcal{H}\!\mathit{om}}\nolimits}
\def\SheafExt{\mathop{\mathcal{E}\!\mathit{xt}}\nolimits}
\def\Sch{\mathit{Sch}}
\def\Mor{\mathop{\mathrm{Mor}}\nolimits}
\def\Ob{\mathop{\mathrm{Ob}}\nolimits}
\def\Sh{\mathop{\mathit{Sh}}\nolimits}
\def\NL{\mathop{N\!L}\nolimits}
\def\CH{\mathop{\mathrm{CH}}\nolimits}
\def\proetale{{pro\text{-}\acute{e}tale}}
\def\etale{{\acute{e}tale}}
\def\QCoh{\mathit{QCoh}}
\def\Ker{\mathop{\mathrm{Ker}}}
\def\Im{\mathop{\mathrm{Im}}}
\def\Coker{\mathop{\mathrm{Coker}}}
\def\Coim{\mathop{\mathrm{Coim}}}

% Boxtimes
%
\DeclareMathSymbol{\boxtimes}{\mathbin}{AMSa}{"02}

%
% Macros for moduli stacks/spaces
%
\def\QCohstack{\mathcal{QC}\!\mathit{oh}}
\def\Cohstack{\mathcal{C}\!\mathit{oh}}
\def\Spacesstack{\mathcal{S}\!\mathit{paces}}
\def\Quotfunctor{\mathrm{Quot}}
\def\Hilbfunctor{\mathrm{Hilb}}
\def\Curvesstack{\mathcal{C}\!\mathit{urves}}
\def\Polarizedstack{\mathcal{P}\!\mathit{olarized}}
\def\Complexesstack{\mathcal{C}\!\mathit{omplexes}}
% \Pic is the operator that assigns to X its picard group, usage \Pic(X)
% \Picardstack_{X/B} denotes the Picard stack of X over B
% \Picardfunctor_{X/B} denotes the Picard functor of X over B
\def\Pic{\mathop{\mathrm{Pic}}\nolimits}
\def\Picardstack{\mathcal{P}\!\mathit{ic}}
\def\Picardfunctor{\mathrm{Pic}}
\def\Deformationcategory{\mathcal{D}\!\mathit{ef}}

% --- appended 2026-10-12 (Chief): the project own definitions of the reference/slogan/history
% display environments, copied mechanically from the shared Stacks preamble lines 8-12
% (\usepackage{verbatim} plus the three \newenvironment lines, which the chapter preamble keeps
% verbatim from the project preamble). Nothing added, removed or redefined: these environments
% carry no mathematical content and the project itself typesets them as comments. Without these
% three lines a counted statement whose own text uses them fails with Environment slogan undefined.
\usepackage{verbatim}
\newenvironment{reference}{\comment}{\endcomment}
\newenvironment{slogan}{\comment}{\endcomment}
\newenvironment{history}{\comment}{\endcomment}
\usepackage{cleveref}
\crefname{originallemma}{original lemma}{original lemmas}
\Crefname{originallemma}{Original lemma}{Original lemmas}
\makeatletter
\newlabel{homology-lemma-exact-functor-ext}{{010O}{source}{}{sourcepointer-1}{}}
\newlabel{homology-lemma-exact-functor-ext@cref}{{[originallemma][0][]010O}{[1][0][]source}{}{}{}}
\makeatother
\crefname{originallemma}{original lemma}{original lemmas}
\Crefname{originallemma}{Original lemma}{Original lemmas}
\makeatletter
\newlabel{homology-lemma-six-term-sequence-ext}{{05E2}{source}{}{sourcepointer-2}{}}
\newlabel{homology-lemma-six-term-sequence-ext@cref}{{[originallemma][0][]05E2}{[1][0][]source}{}{}{}}
\makeatother
\crefname{originallemma}{original lemma}{original lemmas}
\Crefname{originallemma}{Original lemma}{Original lemmas}
\makeatletter
\newlabel{lemma-I-power-torsion-presentation}{{24002}{source}{}{sourcepointer-3}{}}
\newlabel{lemma-I-power-torsion-presentation@cref}{{[originallemma][0][]24002}{[1][0][]source}{}{}{}}
\makeatother
\crefname{originallemma}{original lemma}{original lemmas}
\Crefname{originallemma}{Original lemma}{Original lemmas}
\makeatletter
\newlabel{lemma-explicit-ext}{{24413}{source}{}{sourcepointer-4}{}}
\newlabel{lemma-explicit-ext@cref}{{[originallemma][0][]24413}{[1][0][]source}{}{}{}}
\makeatother
\crefname{originallemma}{original lemma}{original lemmas}
\Crefname{originallemma}{Original lemma}{Original lemmas}
\makeatletter
\newlabel{lemma-koszul-homology-annihilated}{{24454}{source}{}{sourcepointer-5}{}}
\newlabel{lemma-koszul-homology-annihilated@cref}{{[originallemma][0][]24454}{[1][0][]source}{}{}{}}
\makeatother
\crefname{originallemma}{original lemma}{original lemmas}
\Crefname{originallemma}{Original lemma}{Original lemmas}
\makeatletter
\newlabel{lemma-naive-Koszul-complex}{{24394}{source}{}{sourcepointer-6}{}}
\newlabel{lemma-naive-Koszul-complex@cref}{{[originallemma][0][]24394}{[1][0][]source}{}{}{}}
\makeatother
\crefname{originallemma}{original lemma}{original lemmas}
\Crefname{originallemma}{Original lemma}{Original lemmas}
\makeatletter
\newlabel{lemma-neighbourhood-equivalence}{{24301}{source}{}{sourcepointer-7}{}}
\newlabel{lemma-neighbourhood-equivalence@cref}{{[originallemma][0][]24301}{[1][0][]source}{}{}{}}
\makeatother
\crefname{originallemma}{original lemma}{original lemmas}
\Crefname{originallemma}{Original lemma}{Original lemmas}
\makeatletter
\newlabel{lemma-neighbourhood-isomorphism}{{24264}{source}{}{sourcepointer-8}{}}
\newlabel{lemma-neighbourhood-isomorphism@cref}{{[originallemma][0][]24264}{[1][0][]source}{}{}{}}
\makeatother
\title{Extensions of modules over the base glue from I-power torsion extensions after a flat covering}
\author{The Stacks Project Authors}
\date{stacks-project:more-algebra.tex:local-snapshot-20261012; GNU FDL 1.2 or later}
\begin{document}
\maketitle
\setcounter{subsection}{1}
\noindent\textit{Editorial note.} All delivered mathematical text is the author's own, copied from the pinned Stacks Project chapter file more-algebra.tex; no statement or proof step was rewritten, reordered or removed.\par\medskip
\noindent\textit{Editorial note.} The chapter file pinned-original-source.tex carries only \textbackslash{}input\{preamble\} on its first line, so no preamble text lives inside the chapter. The retained file sources/stacks-more-algebra/retained-preamble.tex is a mechanical, line-for-line copy of the project's single shared preamble as staged in the snapshot; nothing was added, removed or redefined and no package was added.\par\medskip
\noindent\textit{Editorial note.} The source URL is the source-level Stacks Project URL, as required when more than one card is produced from one source; the per-claim tag is carried by claim\_key.\par\medskip
\noindent\textit{Editorial note.} The author writes 'Given an short exact sequence' and '(details omitted)' once, in the reduction to M'; both are kept verbatim. The ten references inside the counted proof span are recorded as non-load-bearing navigation pointers; they mix cross-chapter tags with same-chapter label lines.\par\medskip
\noindent\textit{Allowed standard prerequisites (editorial).}\begin{itemize}
\item The neighbourhood equivalence for modules annihilated by a power of I. invoked at proof lines 24523-24524 and 24548-24549. Identifies Hom and Ext groups before and after base change for such modules.
\item The presentation of an I-power torsion module by a direct sum of modules R/I\textasciicircum{}n. invoked at proof lines 24519-24520. Is the reduction that makes the case analysis finite.
\item The explicit description of Ext\textasciicircum{}1(R/I\textasciicircum{}n, N) via Koszul homology and the naive Koszul module. invoked at proof lines 24529-24545. Converts the remaining isomorphism statement into two Hom and homology comparisons that can be checked directly.
\item Flatness giving H\_1(N, f\_bullet) tensor\_R S = H\_1(N tensor\_R S, f\_bullet). used at proof lines 24550-24555. Is the computation that closes the Koszul comparison.
\end{itemize}\medskip
\ifdefined\hypertarget\hypertarget{sourcepointer-1}{}\fi
\noindent\textit{Source reference.} Original lemma 010O: focused-seeds/AGP/sources/stacks-snapshot-20261011/tags.txt, source lines 1334--1334. Cross-chapter navigation pointer into the Stacks Project; the standard fact it supplies is applied inside the counted proof span and the argument resting on it stays there.\par\medskip
\ifdefined\hypertarget\hypertarget{sourcepointer-2}{}\fi
\noindent\textit{Source reference.} Original lemma 05E2: focused-seeds/AGP/sources/stacks-snapshot-20261011/tags.txt, source lines 6975--6975. Cross-chapter navigation pointer into the Stacks Project; the standard fact it supplies is applied inside the counted proof span and the argument resting on it stays there.\par\medskip
\ifdefined\hypertarget\hypertarget{sourcepointer-3}{}\fi
\noindent\textit{Source reference.} Original lemma 24002: sources/stacks-more-algebra/pinned-original-source.tex, source lines 24002--24002. Same-chapter navigation pointer; the statement it names is standard for the section and the argument that consumes it stays inside the counted proof span.\par\medskip
\ifdefined\hypertarget\hypertarget{sourcepointer-4}{}\fi
\noindent\textit{Source reference.} Original lemma 24413: sources/stacks-more-algebra/pinned-original-source.tex, source lines 24413--24413. Same-chapter navigation pointer; the statement it names is standard for the section and the argument that consumes it stays inside the counted proof span.\par\medskip
\ifdefined\hypertarget\hypertarget{sourcepointer-5}{}\fi
\noindent\textit{Source reference.} Original lemma 24454: sources/stacks-more-algebra/pinned-original-source.tex, source lines 24454--24454. Same-chapter navigation pointer; the statement it names is standard for the section and the argument that consumes it stays inside the counted proof span.\par\medskip
\ifdefined\hypertarget\hypertarget{sourcepointer-6}{}\fi
\noindent\textit{Source reference.} Original lemma 24394: sources/stacks-more-algebra/pinned-original-source.tex, source lines 24394--24394. Same-chapter navigation pointer; the statement it names is standard for the section and the argument that consumes it stays inside the counted proof span.\par\medskip
\ifdefined\hypertarget\hypertarget{sourcepointer-7}{}\fi
\noindent\textit{Source reference.} Original lemma 24301: sources/stacks-more-algebra/pinned-original-source.tex, source lines 24301--24301. Same-chapter navigation pointer; the statement it names is standard for the section and the argument that consumes it stays inside the counted proof span.\par\medskip
\ifdefined\hypertarget\hypertarget{sourcepointer-8}{}\fi
\noindent\textit{Source reference.} Original lemma 24264: sources/stacks-more-algebra/pinned-original-source.tex, source lines 24264--24264. Same-chapter navigation pointer; the statement it names is standard for the section and the argument that consumes it stays inside the counted proof span.\par\medskip
\begin{lemma}
\label{lemma-neighbourhood-extensions}
Assume $\varphi : R \to S$ is a flat ring map and $I \subset R$ is a
finitely generated ideal such that $R/I \to S/IS$ is an isomorphism.
Let $M$, $N$ be $R$-modules. Assume $M$ is $I$-power torsion.
Given an short exact sequence
$$
0 \to N \otimes_R S \to \tilde E \to M \otimes_R S \to 0
$$
there exists a commutative diagram
$$
\xymatrix{
0 \ar[r] &
N \ar[r] \ar[d] &
E \ar[r] \ar[d] &
M \ar[r] \ar[d] &
0 \\
0 \ar[r] &
N \otimes_R S \ar[r] &
\tilde E \ar[r] &
M \otimes_R S \ar[r] &
0
}
$$
with exact rows.
\end{lemma}

\begin{proof}
As $M$ is $I$-power torsion we see that $M \otimes_R S = M$, see
Lemma \ref{lemma-neighbourhood-isomorphism}.
We will use this identification without further mention.
As $R \to S$ is flat, the base change functor is exact and we
obtain a functorial map of $\Ext$-groups
$$
\Ext_R(M, N)
\longrightarrow
\Ext_S(M \otimes_R S, N \otimes_R S),
$$
see
Homology, Lemma \ref{homology-lemma-exact-functor-ext}.
The claim of the lemma is that this map is surjective when
$M$ is $I$-power torsion. In fact we will show that it is an
isomorphism. By
Lemma \ref{lemma-I-power-torsion-presentation}
we can find a surjection $M' \to M$ with $M'$ a direct sum of
modules of the form $R/I^n$. Using the long exact sequence of
Homology, Lemma \ref{homology-lemma-six-term-sequence-ext}
and using Lemma \ref{lemma-neighbourhood-equivalence}
we see that it suffices to prove the lemma for $M'$.
Using compatibility of $\Ext$ with direct sums (details omitted)
we reduce to the case where $M = R/I^n$ for some $n$.

\medskip\noindent
Let $f_1, \ldots, f_t$ be generators for $I^n$. By
Lemma \ref{lemma-explicit-ext}
we have a commutative diagram
$$
\xymatrix{
0 \ar[r] &
\Ext_R(R/I^n, N) \ar[r] \ar[d] &
H_1(N, f_\bullet) \ar[r] \ar[d] &
\Hom_R(K, N) \ar[d] \\
0 \ar[r] &
\Ext_S(S/I^nS, N \otimes S) \ar[r] &
H_1(N \otimes S, f_\bullet) \ar[r] &
\Hom_S(K \otimes S, N \otimes S)
}
$$
with exact rows where $K$ is as in
Lemma \ref{lemma-naive-Koszul-complex}.
Hence it suffices to prove that the two right vertical arrows are
isomorphisms. Since $K$ is annihilated by $I^n$ we see that
$\Hom_R(K, N) = \Hom_S(K \otimes_R S, N \otimes_R S)$ by
Lemma \ref{lemma-neighbourhood-equivalence}.
As $R \to S$ is flat we have
$H_1(N, f_\bullet) \otimes_R S = H_1(N \otimes_R S, f_\bullet)$.
As $H_1(N, f_\bullet)$ is annihilated by $I^n$, see
Lemma \ref{lemma-koszul-homology-annihilated}
we have $H_1(N, f_\bullet) \otimes_R S = H_1(N, f_\bullet)$ by
Lemma \ref{lemma-neighbourhood-isomorphism}.
\end{proof}

\end{document}
